\documentclass[11pt]{article}
\usepackage[T1]{fontenc}
\usepackage[margin=1in]{geometry}
\usepackage{listings}
\usepackage{xcolor}
\usepackage{hyperref}

\definecolor{codebg}{rgb}{0.95,0.95,0.95}
\definecolor{codekw}{rgb}{0.0,0.0,0.6}
\definecolor{codecomment}{rgb}{0.3,0.5,0.3}

\lstset{
  basicstyle=\ttfamily\small,
  backgroundcolor=\color{codebg},
  keywordstyle=\color{codekw}\bfseries,
  commentstyle=\color{codecomment}\itshape,
  numbers=left,
  numberstyle=\tiny,
  frame=single,
  breaklines=true,
  showstringspaces=false,
  tabsize=2
}

\hypersetup{colorlinks=true,linkcolor=blue,urlcolor=blue}

\title{The \texttt{ftxc} Command-Line Compiler --- User Manual}
\author{FlashTeX Documentation Team}
\date{\today}

\begin{document}
\maketitle

\section{Introduction}

This manual documents \texttt{ftxc}, a small command-line utility that wraps
the FlashTeX rendering engine for batch document conversion. It is intended
for users who are comfortable invoking tools such as \texttt{pdflatex} or
\texttt{make} from a terminal. The tool is distributed as a single static
binary and requires no additional runtime dependencies once installed.

Throughout this manual we use \texttt{monospace} type to denote literal
commands, file names such as \texttt{main.tex} or \texttt{Makefile}, and
configuration keys like \texttt{output.format}. Inline snippets of code are
written with \lstinline|ftxc build --watch| when a full listing would be
excessive.

For the latest release notes see \url{https://example.org/ftxc/changelog}
or browse the project homepage at \href{https://example.org/ftxc}{the ftxc
project page}. Bug reports should be filed at
\url{https://example.org/ftxc/issues} rather than emailed directly.

\section{Installation}

The recommended installation path is to download a prebuilt release archive
and place the \texttt{ftxc} binary somewhere on your \texttt{PATH}, such as
\texttt{/usr/local/bin} or \texttt{\textasciitilde/.local/bin}. Package
maintainers may instead prefer to build from source using the accompanying
\texttt{Makefile}.

A minimal build from source looks like the transcript in
Listing~\ref{lst:build-c}, which invokes the small bootstrap compiler
written in C that produces the final Rust-based toolchain.

\begin{lstlisting}[language=C,caption={Bootstrap build driver (C)},label={lst:build-c}]
#include <stdio.h>
#include <stdlib.h>

/* Very small driver that shells out to the real build system
   and reports a friendly error if the toolchain is missing. */
int main(int argc, char **argv) {
    const char *target = (argc > 1) ? argv[1] : "release";
    char cmd[256];
    int n = snprintf(cmd, sizeof(cmd), "cargo build --%s", target);
    if (n < 0 || n >= (int)sizeof(cmd)) {
        fprintf(stderr, "error: target name too long\n");
        return 1;
    }
    return system(cmd);
}
\end{lstlisting}

Once installed, verify the toolchain with \lstinline|ftxc --version| and
confirm that the reported version matches the release you downloaded.

\section{Configuration}

\texttt{ftxc} reads a project-local configuration file named
\texttt{ftxc.toml}. A typical configuration enables incremental caching and
selects an output format, as shown in the short Python helper script in
Listing~\ref{lst:gen-config}, which some users run to generate a starter
file programmatically instead of writing \texttt{ftxc.toml} by hand.

\begin{lstlisting}[language=Python,caption={Generating a starter \texttt{ftxc.toml}},label={lst:gen-config}]
import tomllib
import pathlib

DEFAULT_CONFIG = """
[build]
incremental = true
output = "pdf"

[watch]
debounce_ms = 150
ignore = [".git", "target"]
"""

def write_default_config(path: pathlib.Path) -> None:
    if path.exists():
        raise FileExistsError(f"{path} already exists")
    path.write_text(DEFAULT_CONFIG.strip() + "\n")
    # Sanity check: make sure what we just wrote actually parses.
    with path.open("rb") as fh:
        tomllib.load(fh)

if __name__ == "__main__":
    write_default_config(pathlib.Path("ftxc.toml"))
    print("wrote ftxc.toml")
\end{lstlisting}

Configuration keys may also be overridden on the command line, for example
\lstinline|ftxc build --set watch.debounce_ms=50|. Command-line overrides
always take precedence over the values found in \texttt{ftxc.toml}.

\section{Everyday Usage}

The typical development loop consists of editing a \texttt{.tex} source
file, letting \texttt{ftxc} rebuild it, and inspecting the resulting PDF.
The following steps describe that loop:

\begin{itemize}
  \item Run \texttt{ftxc build} once to produce an initial \texttt{main.pdf}.
  \item Run \texttt{ftxc watch} to rebuild automatically whenever
        \texttt{main.tex} or any file it \texttt{\textbackslash input}s changes.
  \item Use \texttt{ftxc diagnose} to print a structured report of warnings,
        including overfull and underfull box warnings by page.
  \item Use \texttt{ftxc clean} to remove \texttt{.aux}, \texttt{.log}, and
        \texttt{.out} artifacts from the working directory.
\end{itemize}

A short interactive session, including the raw \texttt{verbatim} output
printed by the tool, is reproduced below exactly as it appears in a
terminal:

\begin{verbatim}
$ ftxc build
[ftxc] reading ftxc.toml
[ftxc] compiling main.tex (pass 1/2)
[ftxc] compiling main.tex (pass 2/2)
[ftxc] wrote main.pdf (3 pages, 0 overfull, 1 underfull)
$ ftxc diagnose --page 2
[ftxc] page 2: underfull \hbox (badness 4021) in paragraph at lines 40--44
\end{verbatim}

\section{Command-Line Flags}

The complete set of top-level flags accepted by \texttt{ftxc} is listed
below. Flags may be combined freely except where noted.

\begin{description}
  \item[\texttt{-{}-watch}] Rebuild automatically on file changes, using the
        debounce interval from \texttt{watch.debounce\_ms}.
  \item[\texttt{-{}-set \emph{key}=\emph{value}}] Override a single
        configuration key for this invocation only; may be repeated.
  \item[\texttt{-{}-format \emph{fmt}}] Select an output format, one of
        \texttt{pdf}, \texttt{dvi}, or \texttt{ps}.
  \item[\texttt{-{}-jobs \emph{n}}] Limit parallel compilation to \emph{n}
        worker threads; defaults to the number of logical CPUs.
  \item[\texttt{-{}-quiet}] Suppress all output except errors; useful when
        \texttt{ftxc} is invoked from another \texttt{Makefile} target.
  \item[\texttt{-{}-version}] Print the compiler version and exit
        immediately without reading \texttt{ftxc.toml}.
\end{description}

\section{Troubleshooting}

If \texttt{ftxc build} exits with a nonzero status, first check that
\texttt{ftxc.toml} is valid by running \lstinline|ftxc config --check|.
Most remaining failures fall into one of two categories: a missing input
file referenced via \texttt{\textbackslash input} or \texttt{\textbackslash
includegraphics}, or a genuine typesetting error reported by the underlying
engine. In the latter case the error message will include a source file and
line number, much like the diagnostics produced by \texttt{pdflatex}
itself.

For further help, consult \href{https://example.org/ftxc/docs}{the online
documentation} or open an issue with a minimal reproduction attached.

\end{document}
